% slides.tex - Apresentações em Beamer
\documentclass[11pt]{beamer}

% Tema e configurações
\usetheme{Madrid}
\usecolortheme{default}
\useinnertheme{circles}
\useoutertheme{infolines}

% Pacotes
\usepackage[utf8]{inputenc}
\usepackage[T1]{fontenc}
\usepackage[brazil]{babel}
\usepackage{amsmath,amssymb}
\usepackage{graphicx}
\usepackage{listings}
\usepackage{xcolor}
\lstset{
  language=Python,
  basicstyle=\footnotesize\ttfamily,
  numbers=left,
  numberstyle=\tiny,
  frame=single,
  backgroundcolor=\color{gray!10},
  breaklines=true,
  captionpos=b,
  showstringspaces=false,
  keywordstyle=\color{blue}\bfseries,
  commentstyle=\color{green!60!black},
  stringstyle=\color{red}
}
\usepackage{hyperref}

% Configurações de numeração de equações
\numberwithin{equation}{section}
\renewcommand{\theequation}{\thesection.\arabic{equation}}

% Configurações de fonte para seções (removido tamanho especial)
\setbeamerfont{section title}{series=\bfseries,family=\sffamily}
\setbeamerfont{subsection title}{family=\sffamily}

% Template para página de seção com menos espaço
\setbeamertemplate{section page}{
  \begin{centering}
    {\usebeamerfont{section title}\usebeamercolor[fg]{section title}\MakeUppercase{\insertsectionhead}}%
  \end{centering}
}

% Template para página de subsection
\setbeamertemplate{subsection page}{
  \begin{centering}
    {\usebeamerfont{subsection title}\usebeamercolor[fg]{subsection title}\insertsubsectionhead}%
  \end{centering}
}

% Informações da apresentação
\title{Guia de Cálculo Numérico}
\subtitle{Slides de Apresentação}
\author{Pedro}
\institute{Universidade}
\date{\today}

\begin{document}

% Título
\begin{frame}
\titlepage
\end{frame}

% Sumário
\begin{frame}{Sumário}
\tableofcontents
\end{frame}

% Introdução
\section{Introdução}
\begin{frame}{Introdução ao Cálculo Numérico}
\begin{itemize}
\item Importância do cálculo numérico
\item Aplicações práticas
\item Métodos fundamentais
\item Análise de erros
\item Implementações em Python
\end{itemize}
\end{frame}

% Fundamentos
\section{Fundamentos}
\begin{frame}{Representação Numérica}
\begin{itemize}
\item Números de ponto flutuante IEEE 754
\item Precisão simples (32 bits) e dupla (64 bits)
\item Erros de arredondamento e truncamento
\item Eps da máquina
\item Cancelamento catastrófico
\end{itemize}
\end{frame}

\begin{frame}{Análise de Erros}
\begin{itemize}
\item Erro absoluto vs relativo
\item Propagação de erros
\item Condição do problema
\item Número de condição
\item Estabilidade numérica
\end{itemize}
\end{frame}

% Seções dos capítulos
\section{Raízes de Equações}
\begin{frame}{Métodos para Raízes}
\begin{itemize}
\item Método da bissecção (intervalo, convergência garantida)
\item Método de Newton-Raphson (derivada, quadrática)
\item Método da secante (aproximação da derivada, superlinear)
\item Comparação: robustez vs velocidade
\end{itemize}
\end{frame}

\begin{frame}{Bissecção - Algoritmo}
\begin{itemize}
\item Requer intervalo $[a,b]$ com $f(a)f(b) < 0$
\item Convergência linear: erro reduz à metade a cada iteração
\item Sempre converge para uma raiz
\item Fácil implementação, mas lento
\end{itemize}
\end{frame}

\begin{frame}{Newton-Raphson}
\begin{itemize}
\item Fórmula iterativa:
\begin{equation}
x_{n+1} = x_n - \frac{f(x_n)}{f'(x_n)}
\end{equation}
\item Convergência quadrática (ordem 2)
\item Requer derivada e boa estimativa inicial
\item Pode falhar se $f'(x) = 0$
\end{itemize}
\end{frame}

\begin{frame}{Código: Método da Bissecção}
\footnotesize
Implementação básica em Python (código omitido para simplificação da apresentação).
\end{frame}

\begin{frame}{Código: Newton-Raphson}
\footnotesize
Implementação básica em Python (código omitido para simplificação da apresentação).
\end{frame}

\section{Sistemas Lineares}
\begin{frame}{Resolução de Sistemas}
\begin{itemize}
\item Eliminação de Gauss (direta, exata)
\item Decomposição LU (fatoração, reutilização)
\item Método de Jacobi (iterativo, paralelizável)
\item Análise de estabilidade e convergência
\end{itemize}
\end{frame}

\begin{frame}{Eliminação de Gauss}
\begin{itemize}
\item Transforma sistema em triangular superior
\item Operações elementares: troca de linhas, multiplicação
\item Pivotamento parcial para estabilidade
\item Complexidade $O(n^3)$
\end{itemize}
\end{frame}

\begin{frame}{Código: Eliminação de Gauss}
\footnotesize
% Código omitido
\end{frame}

\section{Interpolação}
\begin{frame}{Interpolação Polinomial}
\begin{itemize}
\item Lagrange: forma direta, numericamente estável
\item Newton: diferenças divididas, incremental
\item Splines: evita oscilações, melhor para muitos pontos
\item Fenômeno de Runge: interpolação de alto grau
\end{itemize}
\end{frame}

\begin{frame}{Código: Interpolação de Lagrange}
\footnotesize
% Código omitido
\end{frame}

\section{Ajuste de Curvas}
\begin{frame}{Ajuste de Curvas}
\begin{itemize}
\item Método dos mínimos quadrados
\item Regressão linear: $y = ax + b$
\item Regressão polinomial: grau n
\item Métricas: $R^2$, MSE, MAE
\item Detecção de outliers
\end{itemize}
\end{frame}

\begin{frame}{Código: Regressão Linear}
\footnotesize
% Código omitido
\end{frame}

\section{Integração Numérica}
\begin{frame}{Regra do Trapézio}
\begin{itemize}
\item Aproximação trapezoidal:
\begin{equation}
\int_a^b f(x) \, dx \approx \frac{b-a}{2} (f(a) + f(b))
\end{equation}
\item Regra composta para maior precisão
\item Erro proporcional a $(b-a)^3$
\end{itemize}
\end{frame}

\begin{frame}{Código: Regra do Trapézio}
\footnotesize
% Código omitido
\end{frame}

\section{EDOs}
\begin{frame}{Equações Diferenciais}
\begin{itemize}
\item Método de Euler:
\begin{equation}
y_{n+1} = y_n + h f(x_n, y_n)
\end{equation}
\item Método de Runge-Kutta 4ª ordem: alta precisão
\item Sistemas de EDOs: aplicações reais
\item Estabilidade e erro de truncamento
\end{itemize}
\end{frame}

\begin{frame}{Código: Método de Euler}
\footnotesize
% Código omitido
\end{frame}

\begin{frame}{Código: Runge-Kutta 4ª Ordem}
\footnotesize
% Código omitido
\end{frame}

% Conclusão
\section{Conclusão}
\begin{frame}{Conclusão}
\begin{itemize}
\item Revisão dos conceitos
\item Implementações em Python
\item Próximos passos
\end{itemize}
\end{frame}

\end{document}